\section{Policy and harness co-adaptation}\label{sec:C30:04}
Memory and harness updates change the trajectories on which a model learns. Model updates then change how those same memories and harnesses are interpreted. The coupled process is not generally equivalent to optimizing independent components. Let $\theta$ denote policy parameters and $h$ the harness. A cross-component interaction can be measured by
\begin{equation}
\begin{aligned}
 \Delta_{\rm int}={}&J(\theta_1,h_1)-J(\theta_1,h_0)\\
 &-J(\theta_0,h_1)+J(\theta_0,h_0).
\end{aligned}
\end{equation}
This is a difference of differences, not proof of a universal synergistic mechanism. Its interpretation requires the four configurations to be evaluated on the same task population under matched evaluation budgets. For already trained checkpoints it measures compatibility of those checkpoints and harnesses. Attributing a causal benefit to the training procedure additionally requires controlled training interventions and their complete adaptation budgets. An interaction can change sign across environments. The harness may help the old policy precisely because it compensates for a weakness removed by training.

Self-HPO provides a concrete alternating design: the current policy proposes failure-conditioned harness candidates, then reinforcement learning updates the policy under the selected harness; unresolved tasks inform the next search set. The reported study uses one base policy family and two alternating rounds. Its authors evaluate harness-only, policy-only, and joint changes, including cross-pairings of optimized and base components. This supports studying co-adaptation, while the limited rounds and backbone leave sustained recursive improvement unresolved~\cite{root_selfhpo2026}.

MASS alternates multi-agent workflow evolution and supervised internalization of generated trajectories, using the same underlying model in executor, optimizer, and evaluator roles. The reported experiment demonstrates two cycles. Its architecture raises a useful question: does a richer information topology produce training signals that a single execution would not expose? It does not eliminate evaluator dependence, nor does finite internalization establish that the resulting model can validate arbitrary research conclusions~\cite{root_mass2026}.

Its public research-benchmark analysis further separates delivery from quality. In the reported MLR-Bench runs, the higher aggregate score accompanies more delivered papers; conditional means among delivered papers concern different completed task sets. The appendix reports no delegation calls in those executions and notes that statistical conclusions change with the aggregation unit. Consequently, that transfer result does not itself identify delegation as the mechanism or establish consistent quality improvement across research tasks~\cite{root_mass2026}.

The training clock can also be wrong. Terminal-RSI's second revision, released on 25 October, distinguishes consequential environment turns from observation-token positions when assigning delayed credit. A verbose tool output should not automatically imply a longer decision horizon. Its turn-based discount changes the objective relative to token-time discounting; exact chunk continuation preserves both successor value and successor trace only for unchanged retained inputs. Summarization changes that conditioning. Although the revision describes TB-GAE utilization, its limitations still require an isolated estimator comparison, compute-matched alternation, online reward controls, and inherited-versus-reset critic controls before causal component attribution. Deriving an estimator identity is not evidence that a particular evaluated checkpoint used it~\cite{cumulative_c_terminalrsi_v2}.

These examples motivate a compatibility matrix for every co-evolving system. Model, tokenizer, context transform, memory schema, skill bank, harness, tool schema, reward, and verifier versions should be recorded together. Updating one component requires revalidation of its dependents. Otherwise the next training run may condition on procedures that the behavior policy never used, or a deployment policy may retrieve skills whose preconditions no longer hold.

